\subsection{AI and the Concentration of Power in Tech}
\label{sec:concentration-power}
A handful of firms controls the datasets and infrastructure a competitive system needs, and a firm in that position is a gatekeeper whether or not it intends to be.

The obvious countermeasure is antitrust, and it has been tried. Facebook's acquisitions of Instagram and WhatsApp — the company now operates as Meta — did not just remove two competitors; they handed one firm comprehensive behavioral data from billions of users. The Federal Trade Commission sued in December 2020 to unwind both deals \autocite{ftc2020metaplatforms}. After a six-week bench trial the district court held, in November 2025, that the agency had not shown Meta currently holds monopoly power — TikTok and YouTube compete for the same attention, so a market defined without them was too narrow — and the FTC appealed the following January \autocite{ftc2025metaopinion,ftc2026metaappeal}. Whatever the appeal decides, the case is a lesson about the instrument rather than about this defendant: a merger challenge brought a decade after the mergers has to prove a present-day market, and the market had moved. That leaves two remedies which do not depend on winning that argument. One is mandated access: a dataset or an infrastructure layer that functions as a chokepoint is licensed on published terms to any qualified competitor, which reaches the position rather than the conduct. The other is public funding for research outside the dominant firms, which an earlier section found lasts only where the money is standing rather than project-term. Export controls, prohibited-practices regulation and design standards address the narrower case of authoritarian misuse and are taken up elsewhere; the concern here is prior to that one: concentration itself, whatever any single firm does with the position it holds.
